\begin{figure}[tbp]
\centering
\begin{tikzpicture}[
  archbox/.style={draw=black!65, rounded corners=2pt, line width=.45pt,
    text width=2.94cm, minimum height=1.42cm, inner sep=3.5pt,
    align=center, font=\footnotesize},
  archstep/.style={archbox, fill=black!2},
  archresult/.style={archbox, fill=blue!5},
  archarrow/.style={-{Stealth[length=1.7mm,width=1.2mm]},
    draw=black!75, line width=.55pt},
  archheading/.style={anchor=west, font=\small\bfseries},
  archlabel/.style={font=\scriptsize, fill=white, inner sep=2pt}
]
\node[archheading] at (-1.60,1.26)
  {(a) Numerical inequality};
\node[archstep] (archstress) at (0,0)
  {Auxiliary stress\\ Conserved; flux saturation\\
   Lemmas~\ref{lem:str-flux-saturation}, \ref{lem:str-stress-identities}};
\node[archstep] (archbend) at (3.85,0)
  {Bend to a flat end\\ Cubic mass correction\\
   Proposition~\ref{prop:str-af-bending}};
\node[archstep] (archflat) at (7.70,0)
  {AF coupled deformation\\
   $R_{\widehat g}\ge0$\\
   Lemma~\ref{lem:fc-comparison-metric}};
\node[archstep] (archpenrose) at (11.55,0)
  {Riemannian Penrose\\ Flat mass estimate\\
   Theorem~\ref{thm:fc-mass-completion}};
\draw[archarrow] (archstress) -- (archbend);
\draw[archarrow] (archbend) -- (archflat);
\draw[archarrow] (archflat) -- (archpenrose);

\node[archstep] (archah) at (0,-2.50)
  {AH deformation\\
   $(\widehat g,0)$\\ Mass and area control\\
   Proposition~\ref{prop:ah-existence}};
\node[archresult] (archonesign) at (3.85,-2.50)
  {One-sign bound\\ Applied to $(\widehat g,0)$\\
   Theorem~\ref{thm:str-one-sign}};
\node[archresult] (archcomponent) at (7.70,-2.50)
  {Fixed-chart bound for $g$\\
   $p_0(g)\ge (r_A+r_A^3)/2$\\
   Theorem~\ref{thm:ah-component}};
\node[archresult] (archlorentz) at (11.55,-2.50)
  {Spatial Lorentz\\ Balancing\\
   $m_{\AH}\ge (r_A+r_A^3)/2$\\
   Theorem~\ref{thm:ah-invariant}};
\draw[archarrow] (archpenrose.south) -- (11.55,-1.22)
  -- node[archlabel,above] {ordered limits of inequalities}
  (3.85,-1.22) -- (archonesign.north);
\draw[archarrow] (archah) -- (archonesign);
\draw[archarrow] (archonesign) -- (archcomponent);
\draw[archarrow] (archcomponent) -- (archlorentz);

\node[archheading] at (-1.60,-3.80)
  {(b) Equality on the connected horizon subclass};
\node[archstep] (archvariations) at (0,-5.02)
  {Original-data variations\\ At equality in $(g,K)$\\
   Section~\ref{eq-section}};
\node[archstep] (archadjoint) at (3.85,-5.02)
  {Causal adjoint\\ $(u,X)$; horizon data\\
   Theorem~\ref{eq-causal-killing}};
\node[archstep] (archstatic) at (7.70,-5.02)
  {No interior null set\\ Regular static quotient\\
   Theorem~\ref{geo-static-quotient}};
\node[archstep] (archend) at (11.55,-5.02)
  {Weighted end estimates\\ Normalized smooth end\\
   Proposition~\ref{st-end-regularity}};
\draw[archarrow] (archvariations) -- (archadjoint);
\draw[archarrow] (archadjoint) -- (archstatic);
\draw[archarrow] (archstatic) -- (archend);

\node[archstep] (archhorizon) at (3.85,-7.77)
  {Horizon area and gravity\\
   $|S|_h=4\pi r_h^2$\\
   $\kappa=\frac{1+3r_h^2}{2r_h}$\\
   Proposition~\ref{st-base-equality}};
\node[archresult] (archhw) at (7.70,-7.77)
  {Vanishing static deficit\\ Global static model\\
   Theorem~\ref{hw-static-rigidity}};
\node[archresult] (archgraph) at (11.55,-7.77)
  {Original graph\\
   $(g,K)$, $M=m_{\AH}$\\ Full future horizon\\ Conformal infinity\\
   Proposition~\ref{emb-global}};
\draw[archarrow] (archstatic.south) -- (7.70,-6.16)
  -- (3.85,-6.16) -- (archhorizon.north);
\draw[archarrow] (archend.south) -- (11.55,-6.44)
  -- (7.70,-6.44) -- (archhw.north);
\draw[archarrow] (archhorizon) -- (archhw);
\draw[archarrow] (archhw) -- (archgraph);
\end{tikzpicture}
\caption{Dependencies of the two arguments. AF means asymptotically flat
and AH means asymptotically hyperbolic. Arrows indicate input to the
next step. In (a), the one-sign theorem is first proved by the upper
route and then applied to the AH comparison data; their mass covector
need not be timelike. Only the original mass covector is balanced.
In (b), equality and the numerical inequality give variations of the
original data. The normalized end and the inherited horizon area and
surface gravity jointly supply the static rigidity hypotheses.
Numerical area comparisons use full enclosing cuts.
The equality route uses no limit of the deformation metrics.}
\label{fig:proof-architecture}
\end{figure}
